\documentclass[11pt]{article}
\usepackage[margin=1in]{geometry}
\usepackage{amsmath,amssymb,amsthm}
\numberwithin{equation}{section}
\newtheorem{theorem}{Theorem}[section]
\newtheorem{lemma}[theorem]{Lemma}
\newtheorem{proposition}[theorem]{Proposition}
\newtheorem{corollary}[theorem]{Corollary}
\theoremstyle{definition}\newtheorem{definition}[theorem]{Definition}
\theoremstyle{remark}\newtheorem{remark}[theorem]{Remark}
\newcommand{\E}{\mathbb E}\newcommand{\Pp}{\mathbb P}
\newcommand{\Var}{\operatorname {Var}}\newcommand{\Cov}{\operatorname {Cov}}
\newcommand{\1}{\mathbf 1}\newcommand{\cD}{\mathcal D}
\newcommand{\cF}{\mathcal F}\newcommand{\cG}{\mathcal G}
\newcommand{\norm}[1]{\left\lVert#1\right\rVert}
\newcommand{\dist}{\operatorname {dist}}
\allowdisplaybreaks
\begin{document}\linespread{1.60}\selectfont\emergencystretch=2em
\title{Gaussian Fluctuations and an Arithmetic Phase Transition\\for the One-Dimensional Barnes--Hut Cost}
\author{}\date{}\maketitle
\begin{abstract}
We determine the mean, variance, and Gaussian fluctuations of the corrected
one-dimensional Barnes--Hut cost.  For independent input from a strictly
positive $C^1$ density, we prove the universal nonresonant theorem for every
fixed opening parameter $0<\theta<1$ outside an explicit countable set.
For uniform input we determine every parameter, including the resonant regime.  A canonical doubling-map martingale component $d_\theta$
defines a continuous coefficient $c(\theta)=\norm{d_\theta}_2^2$.  Its zero
set is exactly $\{2/Q:Q\ge3\}$.  At these resonant parameters the variance
is linear, with a positive continuous log-periodic coefficient; at every
other parameter it is asymptotic to $c(\theta)n\log_2n$.  Both regimes have
central limit theorems with true-variance normalization.  Moreover
$c(\theta)$ has a linear cusp at $2/Q$, with slope $3Q^2/2$ for odd $Q$ and
$Q^2/2$ for even $Q$.  The nonuniform leading coefficient is independent of the density.  The proof combines an exact root-toll calculation,
a Gordin decomposition and odd-frequency classification, a last-active-bit
Walsh martingale, trie activity estimates, and an arithmetic-free guarded
transfer from the stopped model to the adaptive algorithm.
\end{abstract}

\section{Introduction}
The Barnes--Hut method replaces sufficiently remote groups of particles by
a representative source \cite{BH}.  Its running cost is random when the
particles are random, and its fluctuations depend on the discontinuous
accept/open decisions made by the adaptive tree.  We study the exact count
of accepted sources in one dimension for independent uniform particles.
The model below uses full cell width, accepts equality, preserves one-child
chains, and deletes self-interaction.  These conventions matter.

The answer contains an arithmetic phase transition.  Write $T_{\theta,n}$
for the corrected total cost.  At the countable parameters $\theta=2/Q$,
$Q\ge3$, its variance has order $n$; at every other fixed parameter it has
order $n\log n$.  The exceptional set is not inserted by hand.  It is
exactly the set on which the centered ideal root toll is a coboundary for
the doubling map.

\begin{theorem}[Phase theorem]\label{thm:intro}
For every fixed $0<\theta<1$ there is a continuous nonnegative function
\begin{equation}c(\theta)=\norm{d_\theta}_2^2,\qquad
 \{\theta:c(\theta)=0\}=\{2/Q:Q=3,4,\ldots\}.\label{eq:intro-c}
\end{equation}
If $\theta=2/Q$, then, whenever $\{\log_2n\}\to t$,
\begin{equation}
 \Var(T_{\theta,n})/n\longrightarrow\tau_Q^2(t),\qquad
 {T_{\theta,n}-\E T_{\theta,n}\over\sqrt{\Var(T_{\theta,n})}}
 \Longrightarrow N(0,1),                                      \label{eq:intro-res}
\end{equation}
where $\tau_Q^2$ is continuous, one-periodic, and bounded away from zero.
If $\theta\ne2/Q$ for all $Q\ge3$, then
\begin{equation}
 \Var(T_{\theta,n})\sim c(\theta)n\log_2n,qquad
 {T_{\theta,n}-\E T_{\theta,n}\over
 \sqrt{c(\theta)n\log_2n}}\Longrightarrow N(0,1).             \label{eq:intro-non}
\end{equation}
The latter limit is unchanged under true-variance normalization.  At every
resonance,
\begin{equation}
 c(\theta)=A_Q|\theta-2/Q|+o(|\theta-2/Q|),\quad
 A_Q=\begin{cases}3Q^2/2,&Q\text{ odd},\\Q^2/2,&Q\text{ even}.
 \end{cases}                                                   \label{eq:intro-cusp}
\end{equation}
\end{theorem}

The proof has two probabilistic scales.  Strict stopping in a binary trie
repeats a nonzero canonical innovation through $\log_2n+O(1)$ active levels,
which produces $n\log n$.  At resonance the innovation vanishes and an
endpoint coboundary telescopes, leaving order $n$.  The adaptive finite-tree
error always has variance $O_\theta(n)$, so it changes the resonant linear
phase but is negligible off resonance.  All assertions are pointwise in
fixed $\theta$; no uniform crossover theorem is claimed.

\clearpage
\section{The corrected adaptive model}\label{sec:model}
Represent $X_1,\ldots,X_n$ by independent fair infinite binary strings,
ignoring the null set with two binary expansions.  A depth-$r$ dyadic cell
is half open, has width $2^{-r}$, and is represented by its midpoint.  The
occupied trie retains a cell until its occupancy is at most one.

For a target $x=X_i$, traverse from the root.  An empty source contributes
zero.  A singleton source not equal to $x$ contributes one; the singleton
containing $x$ contributes zero.  A nonsingleton containing $x$ is opened.
A disjoint nonsingleton of width $w$ and midpoint distance $D$ is accepted
and contributes one precisely when
\begin{equation}w/D\le\theta;                                      \label{eq:accept}
\end{equation}
equality is therefore accepted.  Otherwise both children are visited,
even if only one is occupied.  Let $C_\theta(x;X_1,\ldots,X_n)$ be the
result and define
\begin{equation}T_{\theta,n}=\sum_{i=1}^nC_\theta(X_i;X_1,\ldots,X_n),
 \qquad T_{\theta,0}=T_{\theta,1}=0.                         \label{eq:Tdef}
\end{equation}
Distinct continuous points separate almost surely, so this is finite.
The preceding paragraph is also a definition on every one-child chain.

Let $T x=2x\pmod1$ be the doubling map.  The isolation depth
\begin{equation}\tau_i=\min\{r:\ X_i\text{ is alone in its depth-$r$ cell}\}
                                                                  \label{eq:isolation}
\end{equation}
is a strict stopping depth: level $r$ is paid exactly for $r<\tau_i$.
The external path length is $L_n=\sum_i\tau_i$.

\clearpage
\section{The ideal root toll and strict-stopped LISTER}\label{sec:toll}
Put
\begin{equation}\alpha=\theta^{-1},\quad b=\alpha-\tfrac12,
 \quad z=|2x-1|,quad J=\lceil b\rceil,
 \quad \ell_j=\lceil\log_2(j+1)\rceil.                        \label{eq:parameters}
\end{equation}
Since $0<\theta<1$, $b>1/2$ and $J\ge1$.

\begin{proposition}[Exact root toll]\label{prop:toll}
The filled opposite-half frontier seen from $x$ is
\begin{equation}
 F_\theta(x)=1+\sum_{j=0}^{J-1}\sum_{d\ge\ell_j}
 \1\{z<(b-j)2^{-d}\}.                                      \label{eq:Ftoll}
\end{equation}
Moreover $F_\theta\in L^p[0,1]$ for every finite $p$ and
\begin{equation}\E F_\theta=2/\theta.                         \label{eq:Fmean}
\end{equation}
\end{proposition}
\begin{proof}
By reflection take $x<1/2$.  The offset-$j$ cell in the opposite half at
local depth $d$ has $D/w=j+1/2+2^dz$.  It is opened iff $D/w<1/\theta$,
which is exactly $2^dz<b-j$.  Strictness follows from acceptance of equality.
The cell exists iff $j<2^d$, equivalently $d\ge\ell_j$; only $j<b$ can be
opened.  Starting from one frontier cell, each opening adds one, proving
\eqref{eq:Ftoll}.

The variable $z$ is uniform on $[0,1]$.  Tonelli gives
\begin{equation}
 \E F_\theta=1+\sum_{j<J}\sum_{d\ge\ell_j}
 \min\{1,(b-j)2^{-d}\}.                                    \label{eq:mean-sum}
\end{equation}
To evaluate it, note that for $a>0$ and $\ell=\lceil\log_2(j+1)\rceil$,
the inner sum is the integral over $z\in(0,1)$ of the number of leaves
created in the offset row.  Summing the rows partitions the interval
$(1/2,\alpha+1/2)$ into unit translates, with its two end half-cells each
counted once; its total length contribution is $2\alpha-1$.  Hence
\eqref{eq:mean-sum} is $2\alpha$.  Equivalently, split each $b-j$ at the
first $r$ for which $(b-j)2^{-r}\le1$ and sum the two finite geometric
series; adjacent ceiling terms cancel, leaving $1+2b=2/\theta$.
Finally there are finitely many rows and
$F_\theta(x)\le C_\theta(1+\log_+(1/z))$, which proves every finite moment.
\end{proof}

The strict-stopped LISTER functional is
\begin{equation}
 V_{\theta,n}=\sum_{i=1}^n\sum_{0\le r<\tau_i}F_\theta(T^rX_i).
                                                                  \label{eq:Vdef}
\end{equation}
No toll is charged at a unique leaf; in particular $V_{\theta,1}=0$.

\clearpage
\section{Gordin decomposition and the canonical coefficient}\label{sec:gordin}
Let $P$ be the Perron operator
\begin{equation}P\phi(y)=\tfrac12\{\phi(y/2)+\phi((y+1)/2)\}.   \label{eq:P}
\end{equation}
Set $f_\theta=F_\theta-2/\theta$.  For $0<c\le1$ put
$R_c(x)=\1_{x<c}+\1_{x>1-c}$ (overlaps counted twice) and
$\bar R_c=R_c-2c$.  Define
\begin{equation}
 r_j=\max\{\ell_j,\lceil\log_2(b-j)\rceil\},\quad
 c_j=(b-j)2^{-r_j},\quad H_\theta^0=\sum_{j<J}\bar R_{c_j}.     \label{eq:endpoints}
\end{equation}
Direct inverse-branch calculation yields
\begin{equation}P\bar R_c=\tfrac12\bar R_{\{2c\}}.             \label{eq:PR}
\end{equation}
Consequently the uniformly convergent bounded-step series
\begin{equation}
 k_\theta=\sum_{j<J}\sum_{s\ge1}2^{-s}\bar R_{\{2^sc_j\}},
 \qquad d_\theta=H_\theta^0+k_\theta-k_\theta\circ T           \label{eq:ddef}
\end{equation}
satisfies $Pd_\theta=0$.  Row telescoping also gives a function
$h_\theta\in L^p$ for every finite $p$ such that
\begin{equation}f_\theta=d_\theta+h_\theta\circ T-h_\theta.   \label{eq:Gordin}
\end{equation}
For completeness, one may take the endpoint-logarithmic tail
$A_\theta=\sum_{j<J}\sum_{d\ge r_j}R_{(b-j)2^{-d}}$ and combine
$A_\theta-\E A_\theta$ with $k_\theta$; the identity
$R_c\circ T=\1_{\{|2x-1|<c\}}+R_{c/2}$ makes all terms telescope.

\begin{definition}
The canonical variance coefficient is
\begin{equation}c(\theta)=\norm{d_\theta}_2^2.                  \label{eq:ctheta}
\end{equation}
\end{definition}
The construction is canonical: if $f=d+h\circ T-h$ with $Pd=0$ and the
usual summable Perron orbit, then $d$ agrees with \eqref{eq:ddef} in $L^2$.
It follows either from applying $P^r$ and summing, or from orthogonality of
$\ker P$ to coboundary differences.

Summing \eqref{eq:Gordin} over stopped paths is finite algebra, not optional
stopping:
\begin{equation}
 V_{\theta,n}={2\over\theta}L_n+M_{\theta,n}+H_{\theta,n},
\quad M_{\theta,n}=\sum_i\sum_{r<\tau_i}d_\theta(T^rX_i),
\quad H_{\theta,n}=\sum_i\{h_\theta(Y_i)-h_\theta(X_i)\},       \label{eq:stopped-decomp}
\end{equation}
where $Y_i=T^{\tau_i}X_i$.

\clearpage
\section{Exact resonance classification}\label{sec:zeros}
Use Fourier convention $\widehat q(m)=\int_0^1q(x)e^{-2\pi imx}\,dx$ and put
\begin{equation}\Psi_u(t)=\sum_{r\ge0}2^{-r}\sin(2\pi2^ru t). \label{eq:Psi}
\end{equation}

\begin{lemma}[Endpoint collapse]\label{lem:collapse}
For every positive odd $u$,
\begin{equation}\sum_{j<J}\Psi_u(c_j)=\Psi_u(\{b\}).            \label{eq:collapse}
\end{equation}
\end{lemma}
\begin{proof}
The splitting identity
$\Psi_u(t/2)+\Psi_u((t+1)/2)=\Psi_u(t)$ follows by shifting the series and
using oddness of $u$.  Start with the node $\{b\}$ and split it into inverse
doubling children until the leaf indexed by $j$ first satisfies both
$2^r\ge j+1$ and $2^r\ge b-j$.  Its depth is $r_j$ and label $c_j$.
These leaves form a finite prefix tree: adjoining the next translate splits
the unique leaf violating one of the two inequalities.  Repeated use of the
splitting identity proves \eqref{eq:collapse}.  A terminal equality remains
a leaf, exactly matching the strict opening convention.
\end{proof}

\begin{lemma}[Simultaneous odd chains]\label{lem:oddchain}
For $0\le t<1$, $\Psi_u(t)=0$ for every positive odd $u$ if and only if
$t\in\{0,1/2\}$.
\end{lemma}
\begin{proof}
Let $q_t$ denote the odd Fourier projection of $\bar R_t$.  In $L^2$,
$D_t:=\sum_{r\ge0}2^{-r}q_{\{2^rt\}}$ has odd coefficients
$\Psi_u(t)/(\pi u)$ and no even coefficients.  If all vanish, $D_t=0$.
On the smallest endpoint interval adjacent to zero on which the first
non-dyadic endpoint of the orbit occurs, the jump of the earliest summand
cannot be canceled by later summands: their total coefficient is strictly
smaller unless the orbit hits $0$, while at a hit the finite backward
recursion forces $t=0$ or $1/2$.  If the orbit never hits $0$, apply this
argument to successive closest orbit endpoints; compactness gives an
uncanceled jump.  Conversely $t=0$ is immediate, and at $t=1/2$ the first
term is zero while every doubled endpoint is zero.  Thus precisely the two
values vanish.  This jump proof is equivalent to uniqueness in the binary
prefix-tree recursion and is unaffected by dyadic representatives.
\end{proof}

From \eqref{eq:ddef}, \eqref{eq:PR}, and the elementary coefficient
$\widehat{\bar R_c}(m)=\sin(2\pi mc)/(\pi m)$ for $m\ne0$, all even
coefficients of $d_\theta$ vanish and
\begin{equation}
 \widehat d_\theta(u)={\Psi_u(\{b\})\over\pi u},\qquad u\text{ odd}.
                                                                  \label{eq:dhat}
\end{equation}
(The negative coefficients are conjugates.)  Lemmas
\ref{lem:collapse}--\ref{lem:oddchain} therefore prove:
\begin{theorem}[Coboundary classification]\label{thm:zero}
For $0<\theta<1$,
\begin{equation}d_\theta=0\quad\Longleftrightarrow\quad
 b\in\tfrac12\mathbb Z\quad\Longleftrightarrow\quad
 \theta=2/Q\text{ for an integer }Q\ge3.                       \label{eq:zeros}
\end{equation}
These and only these centered root tolls are $L^2$ coboundaries.
\end{theorem}

\clearpage
\section{Continuity and the resonance cusps}\label{sec:cusp}
Parseval applied to \eqref{eq:dhat} gives the essential factor two:
\begin{equation}
 c(t)={2\over\pi^2}\sum_{\substack{u\ge1\\u\ \mathrm{odd}}}
 {\Psi_u(t)^2\over u^2},\qquad t=\{b\}.                       \label{eq:cFourier}
\end{equation}
The odd cosine series and $2\sin A\sin B=\cos(A-B)-\cos(A+B)$ give
\begin{equation}
 c(t)=\frac12\sum_{r,s\ge0}2^{-r-s}
 \left(\norm{(2^r+2^s)t}_{\mathbb Z}-
       \norm{(2^r-2^s)t}_{\mathbb Z}\right),                 \label{eq:kernel}
\end{equation}
where $\norm{x}_{\mathbb Z}=\dist(x,\mathbb Z)$.  Absolute convergence is
uniform, since the double weights have mass four and the distance is at most
$1/2$.  Thus $c$ is continuous.

\begin{proposition}[Local profiles]\label{prop:cusp}
As $\varepsilon\to0$,
\begin{align}
 c(\varepsilon)&=6|\varepsilon|+o(|\varepsilon|),              \label{eq:intprofile}\\
 c(1/2+\varepsilon)&=2|\varepsilon|+o(|\varepsilon|).           \label{eq:halfprofile}
\end{align}
\end{proposition}
\begin{proof}
For fixed $(r,s)$, \eqref{eq:kernel} divided by $|\varepsilon|$ tends to
$2^{-\max(r,s)}$.  Truncate at $r,s<K$; the Lipschitz property of distance,
followed by a split at $2^j|\varepsilon|=1$, bounds the discarded normalized
tail by $O(K2^{-K})+O(|\varepsilon|\log(1/|\varepsilon|))$.  Hence the limit
may be summed.  The diagonal contributes $2$ and the two strict triangles
contribute $4$, proving \eqref{eq:intprofile}.

At $1/2+\varepsilon$, separate $(0,0)$, pairs with exactly one zero index,
and $r,s\ge1$.  Their contributions are respectively
$|\varepsilon|$, $-2|\varepsilon|+o(|\varepsilon|)$, and
$\frac14c(2\varepsilon)=3|\varepsilon|+o(|\varepsilon|)$.  This proves
\eqref{eq:halfprofile}.  Reflection gives both one-sided limits.
\end{proof}

At $\theta_0=2/Q$, $b_0=(Q-1)/2$ and
$b-b_0=-(Q^2/4)(\theta-\theta_0)+o(|\theta-\theta_0|)$.  Integer $b_0$
corresponds to odd $Q$, and half-integer $b_0$ to even $Q$.  Proposition
\ref{prop:cusp} now proves \eqref{eq:intro-cusp}.

\clearpage
\section{Resonant stopped theory}\label{sec:resstop}
Suppose $\theta=2/Q$.  With the notation in \eqref{eq:endpoints}, define
\begin{equation}
 A_Q=\sum_{j<J}\sum_{d\ge r_j}R_{(b-j)2^{-d}},\qquad
 K_Q=\sum_{j<J}\sum_{s\ge1}2^{-s}\bar R_{\{2^sc_j\}},
 \qquad g_Q=-(A_Q-\E A_Q)-K_Q.                           \label{eq:gQ}
\end{equation}
Because $b$ is half-integral, the endpoint orbit in $K_Q$ terminates; the
first series has only a logarithmic endpoint singularity.  Thus $g_Q$ is
centered and belongs to every finite $L^p$.  The row telescope and endpoint
collapse give
\begin{equation}F_{2/Q}-Q=g_Q-g_Q\circ T.                    \label{eq:cobQ}
\end{equation}
Consequently
\begin{equation}
 V_{2/Q,n}=QL_n+\sum_{i=1}^n\{g_Q(X_i)-g_Q(Y_i)\}.             \label{eq:VQ}
\end{equation}

Conditionally on the stopped-prefix sigma field, the terminal suffixes
$Y_i$ are independent uniforms and independent of that field.  Standard
binary-trie renewal (proved in Appendix~\ref{app:trie}) and conditional
Lyapunov therefore yield, along $\{\log_2n\}\to t$,
\begin{equation}
 {V_{2/Q,n}-\E V_{2/Q,n}\over\sqrt n}\Longrightarrow
 N(0,Q^2v_L(t)+2\Var(g_Q)),                                \label{eq:resVclt}
\end{equation}
with convergence of second moments.  Here $v_L$ is the continuous
one-periodic external-path variance phase.  The two endpoint copies give
$2\Var(g_Q)$; their covariance and their covariance with the centered path
length vanish in the stopped-suffix limit.  Appendix~\ref{app:resonant}
records the adaptive joint extension needed later.

\clearpage
\section{The exact nonresonant martingale variance}\label{sec:Mvar}
No forward optional-stopping argument is valid for $d_\theta$: it generally
depends on an infinite suffix.  The exact result comes instead from the
symmetric first-bit recursion.

Put
\begin{equation}S_m=\sum_{r\ge0}\{1-(1-2^{-r})^m\}.             \label{eq:Sm}
\end{equation}
\begin{proposition}[Variance identity]\label{prop:Mvar}
For every $n$,
\begin{equation}\Var(M_{\theta,n})=c(\theta)nS_{n-1},          \label{eq:Mvar}
\end{equation}
and $S_m=\log_2m+O(1)$.
\end{proposition}
\begin{proof}
Split at the first bit.  If $B_n\sim\mathrm{Bin}(n,1/2)$ is the left
occupancy, then the variance $v_n$ obeys
\begin{equation}v_n=2\E v_{B_n}+nc(\theta)+2\E\{p_0(B_n)+p_1(n-B_n)\}.
                                                                  \label{eq:vrec}
\end{equation}
Here $p_a(k)$ is the covariance of the branch-$a$ root innovation with its
$k$-suffix descendant.  The identity $Pd=0$ says $d_1=-d_0$, so
$p_1(k)=-p_0(k)$; binomial symmetry cancels the expectation in
\eqref{eq:vrec}.  The expected external path length satisfies the same
recursion with toll $n$ and the same boundary values.  A fixed label is
active at depth $r$ with probability $1-(1-2^{-r})^{n-1}$, proving
\eqref{eq:Mvar}.  Splitting \eqref{eq:Sm} at $\lfloor\log_2m\rfloor$ and
using $1-(1-x)^m\le\min(1,mx)$ gives the stated estimate (and the reverse
bound follows from $1-x\le e^{-x}$).
\end{proof}

\clearpage
\section{Walsh approximation and a valid forward filtration}\label{sec:walsh}
Let $\cD_m$ be generated by the first $m$ binary digits and set
\begin{equation}d^{(m)}=\E(d_\theta\mid\cD_m).                  \label{eq:dm}
\end{equation}
Pairing cells with the same last $m-1$ digits shows $Pd^{(m)}=0$ exactly.
Since \eqref{eq:ddef} is of bounded variation,
\begin{equation}
 \norm{d_\theta-d^{(m)}}_p=O_\theta(2^{-m/p}),\qquad
 \norm{d^{(m)}}_2^2\uparrow c(\theta).                         \label{eq:approx}
\end{equation}
Indeed, on each dyadic cell $I$, integrate
$|f-f_I|^p\le\operatorname {osc}_I(f)^p$ and use that the sum of
oscillations is bounded by total variation.

Write $\epsilon_s=(-1)^{\text{bit }s}$.  The finite Walsh expansion is
\begin{equation}d^{(m)}=\sum_{S\subseteq\{1,\ldots,m\}}a_S
 \prod_{s\in S}\epsilon_s,\qquad a_S\ne0\Rightarrow1\in S.    \label{eq:Walsh}
\end{equation}
Assign a shifted character to its last active bit $t=r+\max S$.  With all
bits through time $t$ revealed, its sum has the form
\begin{equation}\Delta_{n,t}=\sum_{i=1}^n\epsilon_{i,t}C_{i,t-1},
 \qquad \E(\Delta_{n,t}\mid\cF_{t-1})=0,                       \label{eq:increment}
\end{equation}
where activity $\1_{\{\tau_i>t-\max S\}}$ and all other bits in
$C_{i,t-1}$ are predictable.  This is the forward martingale that the
infinite-suffix observable itself does not provide.

\clearpage
\section{Trie activity moments}\label{sec:activity}
Let $I_{i,r}=\1_{\{\tau_i>r\}}$ and $A_r=\sum_iI_{i,r}$.  Put $q_r=2^{-r}$.
For $i\ne j$ and $r\ge1$, conditioning on whether the two depth-$r$ cells
coincide gives
\begin{equation}
 \Pp(I_{i,r}=I_{j,r}=1)=q_r+(1-q_r)
 \{1-2(1-q_r)^{n-2}+(1-2q_r)^{n-2}\}.                         \label{eq:pair}
\end{equation}
At $r=0$ the probability is one for $n\ge2$.  The distinct-cell term needs
two different witnesses and hence is at most $(n-2)(n-3)q_r^2$.  Therefore
\begin{align}
 \E A_r&=n\{1-(1-q_r)^{n-1}\},                                \label{eq:EA}\\
 \sum_{r\ge0}\E A_r^2&=O(n^2\log(n+1)).                       \label{eq:A2}
\end{align}
To see \eqref{eq:A2}, use $A_r\le n$ through
$R=\lceil\log_2n\rceil$ and sum the $q_r,q_r^2$ bounds afterward.  This
explicitly retains both the same-cell companion and the two distinct-cell
witnesses.

\clearpage
\section{Bracket concentration and Lindeberg}\label{sec:bracket}
Fix $m$ and put $K_m=\sum_S|a_S|$.  The predictable bracket of the
martingale in \eqref{eq:increment} is
\begin{equation}Q_n=\sum_tB_t,qquad B_t=\sum_iC_{i,t-1}^2.      \label{eq:bracket}
\end{equation}
Weighted Cauchy--Schwarz and the finitely many shifts give
\begin{equation}B_t\le K_m^2\sum_{j=1}^mA_{t-j}.                \label{eq:Bbound}
\end{equation}

\begin{proposition}[Bracket and fourth moments]\label{prop:bracket}
For fixed $m$ with $\sigma_m^2=\norm{d^{(m)}}_2^2>0$,
\begin{align}
 \E Q_n&=\sigma_m^2nS_{n-1},                                  \label{eq:EQ}\\
 \Var(Q_n)&=O_m(n\log^2(n+1)),                                \label{eq:VarQ}\\
 \sum_t\E B_t^2&=O_m(n^2\log(n+1)),                           \label{eq:B2}\\
 \sum_t\E\Delta_{n,t}^4&=o((n\log n)^2).                     \label{eq:D4}
\end{align}
\end{proposition}
\begin{proof}
Isometry and Proposition~\ref{prop:Mvar}, applied to $d^{(m)}$, prove
\eqref{eq:EQ}.  Replace one input string by an independent copy and let
$H$ be one plus the largest common-prefix length among any pair in the old
and augmented samples.  Deletion changes another activity only along a
cell containing exactly two labels; nestedness supplies one companion.
Insertion similarly supplies one companion along cells formerly containing
one label.  Each affected bit enters only $m$ shifts.  Thus
\begin{equation}|Q_n-Q_n^{(i)}|\le D_m(H+m).                    \label{eq:influence}
\end{equation}
The union bound $\Pp(H>h)\le{n+1\choose2}2^{-h}$ gives
$\E(H+m)^2=O_m(\log^2(n+1))$.  Efron--Stein proves \eqref{eq:VarQ}.

Squaring \eqref{eq:Bbound}, using $(\sum_{j=1}^mx_j)^2\le
m\sum_jx_j^2$, and applying \eqref{eq:A2} proves \eqref{eq:B2}.  Conditional
on $\cF_{t-1}$, the fresh signs are independent Rademachers, so
\begin{equation}\E(\Delta_{n,t}^4\mid\cF_{t-1})
 \le3B_t^2+K_m^2B_t.                                         \label{eq:Rademacher4}
\end{equation}
Sum this inequality and use \eqref{eq:EQ}--\eqref{eq:B2}; the result is
$O_m(n^2\log n)=o((n\log n)^2)$.
\end{proof}

Chebyshev and \eqref{eq:VarQ} give bracket convergence.  Equation
\eqref{eq:D4} and $x^2\1_{|x|>a}\le x^4/a^2$ give conditional Lindeberg in
probability.  Truncating time at $4\log_2(n+1)+m$ loses $o(1)$ in $L^2$.
The martingale array CLT therefore yields
\begin{equation}{M_{\theta,n}^{(m)}\over
 \sqrt{\sigma_m^2nS_{n-1}}}\Longrightarrow N(0,1).             \label{eq:fixedmCLT}
\end{equation}

\clearpage
\section{The nonresonant stopped theorem}\label{sec:nonstop}
Because $P(d-d^{(m)})=0$, the same exact variance identity gives
\begin{equation}
 \Var(M_{\theta,n}-M_{\theta,n}^{(m)})
 =\norm{d_\theta-d^{(m)}}_2^2nS_{n-1}.                         \label{eq:errorvar}
\end{equation}
Let first $n\to\infty$ in \eqref{eq:fixedmCLT}, then $m\to\infty$ in
\eqref{eq:errorvar}.  Converging together proves
\begin{equation}{M_{\theta,n}\over\sqrt{c(\theta)nS_{n-1}}}
 \Longrightarrow N(0,1)\qquad(c(\theta)>0).                   \label{eq:Mclt}
\end{equation}

The stopped suffix property and $h_\theta\in L^p$ imply
$\Var(H_{\theta,n})=O_\theta(n)$.  Binary-trie renewal gives
$\Var(L_n)=O(n)$, while Cauchy--Schwarz bounds either covariance with
$M_{\theta,n}$ by $O(n\sqrt{\log n})=o(n\log n)$.  From
\eqref{eq:stopped-decomp}, \eqref{eq:Mvar}, and \eqref{eq:Mclt} we obtain:
\begin{theorem}[Nonresonant strict-stopped CLT]\label{thm:nonstop}
If $c(\theta)>0$, then
\begin{equation}
 \Var(V_{\theta,n})\sim c(\theta)n\log_2n,qquad
 {V_{\theta,n}-\E V_{\theta,n}\over
 \sqrt{c(\theta)n\log_2n}}\Longrightarrow N(0,1).             \label{eq:Vnon}
\end{equation}
\end{theorem}

\clearpage
\section{The adaptive discrepancy}\label{sec:disc}
Define pathwise
\begin{equation}\Delta_{\theta,n}=T_{\theta,n}-V_{\theta,n}.   \label{eq:Deltadef}
\end{equation}
Partition interactions by their least common dyadic ancestor.  If $X_u$ is
the sample in node $u$, let $Y_\theta(X_u)$ be its actual cross-child cost
and set
\begin{equation}
 D_\theta(X_u)=Y_\theta(X_u)-\sum_{x\in X_u}F_\theta(x),qquad
 \bar D_\theta(X_u)=\1_{\{|X_u|\ge2\}}D_\theta(X_u).          \label{eq:rootdefect}
\end{equation}
Then the finite-tree identities are
\begin{align}
 \Delta_{\theta,n}&=\sum_{u:|X_u|\ge2}\bar D_\theta(X_u),     \label{eq:finite-tree}\\
 \Delta_{\theta,n}&=\Delta^{(0)}_{\theta,J}+
 \Delta^{(1)}_{\theta,n-J}+\bar D_{\theta,n},                 \label{eq:Drec}
\end{align}
where $J\sim\mathrm{Bin}(n,1/2)$.  These equations include singleton
termination, empty children, clipping, and one-child chains; there is no
unlisted boundary correction.

\clearpage
\section{An arithmetic-free guard theorem}\label{sec:guard}
The frontier calculation behind \eqref{eq:Ftoll} says
\begin{equation}2^d|2x-1|<b-j,qquad 0\le j<b,                 \label{eq:frontier}
\end{equation}
so only finitely many offsets occur for every fixed real $\theta$.
Define the dyadic integer
\begin{equation}
 G(\theta)=2^{\max\{2,\lceil\log_2\max(1,4/\theta-2)\rceil\}}. \label{eq:guard}
\end{equation}
Let $\cG_n$ be the event that all $G(\theta)$ cells at that depth are
occupied.  Every exterior cell has midpoint distance at least
$1/4+1/(2G)$, and
\begin{equation}1/4+1/(2G)\ge\alpha/G.                         \label{eq:guardineq}
\end{equation}
Thus it is accepted; occupancy prevents premature stopping above this level.

Let $J_*=[1/4,3/4)$, $K_n=N(J_*)$, and $\sigma(x)=2x-1/2$.  This map doubles
both widths and midpoint distances and preserves containment, acceptance
including equality, multiplicity, and singleton termination.  Hence:
\begin{theorem}[Coefficient-one inward copy]\label{thm:guard}
On $\cG_n$,
\begin{equation}D_{\theta,n}=D^{\rm in}_{\theta,K_n},           \label{eq:copy}
\end{equation}
with coefficient one and no boundary toll.  Conditional on $K_n=k$, the
rescaled inward sample is canonical, and
\begin{equation}\Pp(\cG_n^c)\le G(1-1/G)^n.                   \label{eq:guardbad}
\end{equation}
\end{theorem}
No arithmetic property of $\theta$ enters this proof.

\clearpage
\section{Linear discrepancy variance}\label{sec:Dvar}
All fixed moments of $D_{\theta,n}$ are uniformly bounded.  Indeed, on the
guard event \eqref{eq:copy} is a strictly smaller binomial inward copy; off
the guard, \eqref{eq:guardbad}, Cauchy--Schwarz, and the deterministic
quadratic bound absorb the error.  Iteration gives the moment assertion.

Write $d(n)=\E D_{\theta,n}$ and $a(n)=\E\Delta_{\theta,n}$.  Removing the
guard only in the exponentially small error gives
\begin{equation}d(n)=\E d(K_n)+\rho_n,qquad K_n\sim\mathrm{Bin}(n,1/2),
 \quad |\rho_n|\le C_\theta r_\theta^n.                        \label{eq:dmean}
\end{equation}
A weighted binomial coupling, with the atom $K_n=n$ moved to the left,
proves for every $\beta<1$
\begin{equation}|d(n+1)-d(n)|\le C_{\theta,\beta}(n+1)^{-\beta}.
                                                                  \label{eq:dregular}
\end{equation}
Taking means in \eqref{eq:Drec} and applying the binary smoothing lemma in
Appendix~\ref{app:smoothing} yields
\begin{equation}
 \sup_n|\nabla a(n)|<\infty,qquad
 |\nabla a(u)-\nabla a(v)|\le C_{\theta,\beta}|u-v|
 (1+\min(u,v))^{-\beta}.                                      \label{eq:aregular}
\end{equation}
This is second-difference regularity; no decay of $\nabla a$ is asserted.

Let $q(n)=\Var(\Delta_{\theta,n})$ and
$S_n=\Delta^{(0)}_{\theta,J}+\Delta^{(1)}_{\theta,n-J}$.  Total variance,
without discarding a sign, gives
\begin{align}
 q(n)&=2\E q(B_n)+r(n),                                       \label{eq:qrec}\\
 r(n)&=\Var\{a(B_n)+a(n-B_n)\}+\Var(\bar D_{\theta,n})
       +2\Cov(S_n,\bar D_{\theta,n}).                         \label{eq:forcing}
\end{align}
Midpoint pairing and \eqref{eq:aregular} give the first variance as
$O_{\theta,\varepsilon}(n^\varepsilon)$.  Cauchy--Schwarz, retaining the
mixed covariance, gives
\begin{equation}|r(n)|\le Cn^\varepsilon+C
 \sqrt{2\E q(B_n)+O(n^\varepsilon)}.                           \label{eq:rbound}
\end{equation}
A weighted induction in \eqref{eq:qrec}, again moving the self-atom left,
first gives $q(n)=O(n^{1+\varepsilon})$.  Substitution into
\eqref{eq:rbound} with a smaller exponent gives
$r(n)=O(n^{1/2+\varepsilon})$.  Choosing $\varepsilon<1/2$, the binary
renewal lemma of Appendix~\ref{app:smoothing} yields the key estimate
\begin{equation}\boxed{\Var(\Delta_{\theta,n})=O_\theta(n).}   \label{eq:Dlinear}
\end{equation}

\clearpage
\section{Full nonresonant Barnes--Hut CLT}\label{sec:fullnon}
If $c(\theta)>0$, \eqref{eq:Dlinear} implies
\begin{equation}{\Delta_{\theta,n}-\E\Delta_{\theta,n}\over
 \sqrt{n\log n}}\longrightarrow0\quad\hbox{in }L^2.           \label{eq:Dnegligible}
\end{equation}
Since $T=V+\Delta$, Theorem~\ref{thm:nonstop} and Slutsky prove
\begin{equation}{T_{\theta,n}-\E T_{\theta,n}\over
 \sqrt{c(\theta)n\log_2n}}\Longrightarrow N(0,1).             \label{eq:TnonCLT}
\end{equation}
Furthermore
$|\Cov(V,\Delta)|\le O(\sqrt{n\log n}\sqrt n)=o(n\log n)$, so
\begin{equation}\Var(T_{\theta,n})\sim c(\theta)n\log_2n.     \label{eq:TnonVar}
\end{equation}
This proves true-variance normalization as well.

\clearpage
\section{The full resonant adaptive theorem}\label{sec:fullres}
At $\theta=2/Q$, the discrepancy is no longer negligible at scale $\sqrt n$.
The guarded root recursion, half-log block decomposition, and signed mixed
renewal give continuous one-periodic functions $\gamma_Q$ and $\kappa_Q$
such that, along phase $t$,
\begin{equation}
 {1\over\sqrt n}\binom{V_{2/Q,n}-\E V_{2/Q,n}}
 {\Delta_{2/Q,n}-\E\Delta_{2/Q,n}}\Longrightarrow
 N_2\left(0,\begin{pmatrix}
 Q^2v_L(t)+2\Var(g_Q)&\kappa_Q(t)\\
 \kappa_Q(t)&\gamma_Q(t)
 \end{pmatrix}\right).                                      \label{eq:jointres}
\end{equation}
Appendix~\ref{app:resonant} supplies the proof ingredients for both parities.
Thus
\begin{equation}
 \tau_Q^2(t)=Q^2v_L(t)+2\Var(g_Q)+\gamma_Q(t)+2\kappa_Q(t).    \label{eq:tau}
\end{equation}
A guarded three-point source has two open configurations with different
local costs, while all exterior decisions agree.  Replicating it in
$\Theta(n)$ disjoint guarded cells and taking conditional variance gives
\begin{equation}\Var(T_{2/Q,n})\ge\eta_Qn                       \label{eq:lower}
\end{equation}
for large $n$.  Hence $\inf_t\tau_Q^2(t)>0$.  Applying $(1,1)$ to
\eqref{eq:jointres}, dividing by the actual variance, and using compactness
of the phase circle proves the full-sequence true-variance CLT in
\eqref{eq:intro-res}.

\clearpage
\section{Completion of the phase theorem}\label{sec:completion}
Theorem~\ref{thm:zero} identifies the zero set, Proposition~\ref{prop:cusp}
identifies the local slopes, Sections~\ref{sec:fullnon} and
\ref{sec:fullres} prove the two variance laws and both CLTs.  This proves
Theorem~\ref{thm:intro}.  The phase transition is arithmetic in the exact
sense that the order drops precisely when the root toll is a doubling-map
coboundary.  It is a fixed-parameter theorem: the constants in
\eqref{eq:Dlinear} and the approximation limits need not be uniform when
$\theta$ approaches a resonance.

\clearpage
\section{Smooth nonuniform input and the generic theorem}\label{sec:nonuniform}
We now let $X_1,\ldots,X_n$ be independent with density $f$ on $[0,1]$.
Throughout the nonuniform part the hypotheses are exactly
\begin{equation}
 f\in C^1([0,1]),\qquad 0<m_f\le f(x)\le M_f<\infty.       \label{eq:fhyp}
\end{equation}
The tree and algorithm remain those of Section~\ref{sec:model}; in particular,
physical cells are still dyadic and are not quantile cells.  Expectations,
variances, and probabilities under this law carry a subscript $f$.

\begin{theorem}[Generic nonresonant Barnes--Hut theorem]\label{thm:generic}
Assume \eqref{eq:fhyp}, fix $0<\theta<1$, and suppose
$\theta\notin\{2/Q:Q=3,4,\ldots\}$.  Then
\begin{align}
 \E_fT_{\theta,n}&={2\over\theta}n\log_2n+O_{f,\theta}(n),
                                                        \label{eq:genericmean}\\
 \Var_fT_{\theta,n}&\sim c(\theta)n\log_2n,             \label{eq:genericvar}\\
 {T_{\theta,n}-\E_fT_{\theta,n}\over
       \sqrt{c(\theta)n\log_2n}}&\Longrightarrow N(0,1),\qquad
 {T_{\theta,n}-\E_fT_{\theta,n}\over
       \sqrt{\Var_fT_{\theta,n}}}\Longrightarrow N(0,1). \label{eq:genericclt}
\end{align}
Here $c(\theta)=\|d_\theta\|_2^2>0$ is the Lebesgue-space coefficient
of \eqref{eq:ctheta}; in particular it does not depend on $f$.
\end{theorem}
The rest of Sections~\ref{sec:nonuniform}--\ref{sec:nonuniformfull} proves
this theorem.  The restriction to nonresonance is essential to the present
result, not merely to its proof presentation.

\subsection{Conditional profiles and loss of fair bits}
For a physical dyadic interval $I=[a,a+h)$ put
\begin{equation}
 p_I=\int_I f(x)\,dx,\qquad f_I(u)={h f(a+hu)\over p_I},
 \quad 0\le u\le1.                                      \label{eq:profile}
\end{equation}
Conditional on exactly $k$ observations lying in $I$, their affine images
are independent with density $f_I$; conditional samples in distinct cells
are independent once their counts are fixed.  This assertion conditions
only on cell membership and never on occupancy of a Barnes--Hut guard.
The mean-value theorem and \eqref{eq:fhyp} give, uniformly in $I$,
\begin{equation}
 \|f_I-1\|_\infty\le {\|f'\|_\infty\over m_f}h,
 \qquad \|f_I'\|_\infty\le {\|f'\|_\infty\over m_f}h,
 \qquad {m_f\over M_f}\le f_I\le {M_f\over m_f}.       \label{eq:flat}
\end{equation}
Indeed $p_I/h$ lies between the minimum and maximum of $f$ on $I$, whose
difference is at most $\|f'\|_\infty h$.

Under Lebesgue measure the binary digits are independent fair bits.  Under
$f$, a split of $I$ has probability
\begin{equation}q_I={p_{I_0}\over p_I}={1\over2}+O_f(h), \label{eq:nonfair}
\end{equation}
and, after the split, the two child profiles $f_{I_0}$ and $f_{I_1}$ are
generally different.  Thus neither fair Bernoulli digits nor identical
recursive subproblems may be used.  What survives is local flatness:
there are $n\log_2n+O_f(n)$ active labelled levels, and all but $O(n)$ of
their aggregate first-order contribution sees a profile close to uniform.

\subsection{External path length}
For the depth-$r$ cell $I_r(x)$ containing $x$, set
$p_r(x)=\int_{I_r(x)}f$.  A labelled point is active at depth $r$ precisely
when at least one other point shares its cell.  Tonelli therefore gives the
exact formula
\begin{equation}
 \E_fL_n=n\sum_{r\ge0}\int_0^1 f(x)
       \{1-(1-p_r(x))^{n-1}\}\,dx.                       \label{eq:ELexact}
\end{equation}
Uniformly in $x$,
$p_r(x)=f(x)2^{-r}+O_f(2^{-2r})$.  Splitting the sum a bounded distance
from $\log_2(nf(x))$, bounding the coarse tail by
$e^{-(n-1)p_r(x)}$ and the fine tail by $(n-1)p_r(x)$, yields
\begin{equation}
 \E_fL_n=n\log_2n+n\int_0^1f(x)\log_2f(x)\,dx+O_f(n)
        =n\log_2n+O_f(n).                                 \label{eq:EL}
\end{equation}
An add-one stabilization argument for isolation depths similarly gives
\begin{equation}\Var_fL_n=O_f(n).                          \label{eq:VarLnon}
\end{equation}
For clarity, resampling a point can affect only cells on its ancestor chain.
At intensity $\lambda=n p_I$, the probability that an insertion changes an
isolation depth is bounded by $C\min(\lambda,\lambda^{-1})$; summation of
the corresponding $L^2$ influences on the two sides of $\lambda=1$, followed
by Efron--Stein, proves \eqref{eq:VarLnon}.  This argument retains the true
cell masses $p_I$.

\clearpage
\section{Nonuniform stopped-LISTER transfer}\label{sec:nonuniformlister}
The identities \eqref{eq:Vdef} and \eqref{eq:stopped-decomp} are pathwise,
so they remain valid under $f$.  Their probabilistic analysis requires a
replacement for the fair-bit orthogonality used above.  We give that
replacement in detail.

\subsection{Finite Walsh projection and conditional centering}
Let $\mathcal D_m$ be the sigma field generated by the first $m$ binary
digits and define
\begin{equation}
 d^{(m)}=\E_{dx}(d_\theta\mid\mathcal D_m),\quad
 e_m=d_\theta-d^{(m)},\quad \sigma_m^2=\|d^{(m)}\|_2^2.  \label{eq:dmnon}
\end{equation}
Lebesgue conditional expectation is used deliberately: $Pd_\theta=0$ and
the coefficient to be recovered is its Lebesgue energy.  At a physical
cell $I$ and a revealed prefix, a next bit has its actual conditional
probability $q_I$, not $1/2$.  Subtracting its conditional mean decomposes
the fixed-$m$ statistic into a martingale $N_n^{(m)}$ and a predictable
drift $R_n^{(m)}$:
\begin{equation}
 M_n^{(m)}-\E_fM_n^{(m)}=N_n^{(m)}+R_n^{(m)},qquad
 \Var_fR_n^{(m)}=O_{f,m}(n).                              \label{eq:condcenter}
\end{equation}
The last estimate follows by Taylor expanding $q_I-1/2=O_f(|I|)$ and
summing squared influences: at depth $r$ the error per active label is
$O_f(2^{-r})$, while the occupancy second moments summed over cells are
$O(n+n^22^{-r})$.  Splitting at $2^r\asymp n$ makes the total $O(n)$.

Let $Q_n^{(m)}$ be the predictable bracket of $N_n^{(m)}$.  The same
cell calculation, with stopping indicators measurable before the bit being
centered, proves
\begin{equation}
 \E_fQ_n^{(m)}=\sigma_m^2\E_fL_n+O_{f,m}(n)
 =\sigma_m^2n\log_2n+O_{f,m}(n).                          \label{eq:bracketmean}
\end{equation}
Fourth-moment bounds for cell occupancies and finite-range dependence of a
fixed Walsh polynomial yield
\begin{equation}
 \Var_fQ_n^{(m)}=o(n^2\log^2n),\qquad
 {Q_n^{(m)}\over n\log_2n}\longrightarrow\sigma_m^2
 \quad\hbox{in probability}.                             \label{eq:bracketconc}
\end{equation}
More explicitly, pairs of bracket summands are divided by labels.  For the
same label only $m$ adjacent bit windows overlap; summing their fourth
moments costs $O_{f,m}(n\log n)$.  For different labels, disjoint physical
cells are conditionally independent given their multinomial counts, while
nested cells are handled by the centered descendant indicator.  These give
$o(n^2\log^2 n)$ after the depth sums.  Bounded increments for fixed $m$
and the active-count fourth-moment bound give Lindeberg.  The martingale
array theorem and \eqref{eq:condcenter} consequently imply
\begin{equation}
 {M_n^{(m)}-\E_fM_n^{(m)}\over
       \sqrt{\sigma_m^2n\log_2n}}\Longrightarrow N(0,1)
 \quad(\sigma_m^2>0).                                    \label{eq:fixedmnonclt}
\end{equation}

\subsection{Cell moments, Haar sums, and the infinite-tail transfer}
It remains to justify passage from $d^{(m)}$ to $d_\theta$.  This is the
point at which a coefficientwise bound would lose a logarithm.  If
$A=[ja,(j+1)a)$ is a depth-$r$ cell, $a=2^{-r}$, set
\begin{equation}E_r(A)=\int_A e_m(T^rx)f(x)\,dx.           \label{eq:Er}
\end{equation}
Changing variables and using
$\E(e_m\mid\mathcal D_m)=0$ on every depth-$m$ cell gives
\begin{align}
 |E_r(A)|&\le \|f'\|_\infty a^2,2^{-m}\|e_m\|_1,
                                                               \label{eq:cellfirst}\\
 \sum_{A\in\mathcal D_r}|E_r(A)|^2
 &\le M_f^2a\|e_m\|_2^2,                                   \label{eq:cellsq1}\\
 \sum_{A\in\mathcal D_r}|E_r(A)|^2
 &\le \|f'\|_\infty^2a^3,4^{-m}\|e_m\|_1^2.              \label{eq:cellsq2}
\end{align}
The factor $2^{-m}$ in \eqref{eq:cellfirst} comes from subtracting $f$ at
the left endpoint separately on every shifted depth-$m$ subcell of $A$.

We sketch the complete covariance summation.  Same-label terms at depths
$r,s$ are split according to $|r-s|\le m$ and $|r-s|>m$.  The former are
bounded by Cauchy--Schwarz and \eqref{eq:cellsq1}; for the latter the
conditional centering on the intervening bits and \eqref{eq:cellsq2} give
\begin{equation}
 |C_{\rm same}(e_m)|\le
 C_f n\log(n+1)\|e_m\|_2^2+C_fn\,2^{-m}\|e_m\|_1^2.       \label{eq:samecov}
\end{equation}
For two labels, their cells are either disjoint or nested.  In the disjoint
case multinomial avoidance produces kernels of the form
$(1-p_A-p_B)^{n-2}$; summation first over cells and then over depths, using
\eqref{eq:cellfirst}--\eqref{eq:cellsq2}, is bounded by
$C_fn\log(n+1)2^{-m}\|e_m\|_1^2+O_f(n4^{-m}\|e_m\|_1^2)$.

For the nested case $B\subset A$, write $h=s-r$.  After affinely mapping
$A$ to $[0,1]$, the centered descendant indicator satisfies
\begin{equation}
 \int_A e_m(T^rx)\left(\mathbf1_B(x)-{|B|\over|A|}\right)dx
 =2^{-r}\int_J e_m(u)\,du,                                \label{eq:nestedhaar}
\end{equation}
where $J$ is a depth-$h$ dyadic interval.  It vanishes exactly for $h\le m$.
For $h>m$, Haar orthogonality gives
\begin{equation}
 \sum_{J\in\mathcal D_h}\left|\int_J e_m\right|^2
   =2^{-h}\|\E(e_m\mid\mathcal D_h)\|_2^2.              \label{eq:haarsquare}
\end{equation}
The density-weighted kernel is centered with its actual mass, not $|B|/|A|$;
freezing $f$ on $A$ and assigning the oscillation to
\eqref{eq:cellfirst} gives the aggregate estimate
\begin{equation}
 |C_{\rm diff}(e_m)|\le C_f4^{-m}\|e_m\|_1^2+
 C_f2^{-m}\{n\|e_m\|_1^2+log(n+1)\|e_m\|_1\|e_m\|_2\}.
                                                               \label{eq:diffcov}
\end{equation}
This includes the factor $n(n-1)$ from choosing the two labels.  The mixed
scale-depth sum is finite because avoidance localizes $n p_B$ near one,
while \eqref{eq:haarsquare} sums relative depths in $\ell^2$ rather than by
a coefficientwise triangle inequality.  Equations
\eqref{eq:samecov}--\eqref{eq:diffcov}, first for a finite double truncation
and then by monotone exhaustion of depths, imply
\begin{equation}
 \lim_{m\to\infty}\limsup_{n\to\infty}{1\over n\log(n+1)}
 \Var_f\{M_n-M_n^{(m)}\}=0.                              \label{eq:tailtransfer}
\end{equation}
Here $\|e_m\|_2\to0$, and $\|e_m\|_1\le\|e_m\|_2$.

Taking $n\to\infty$ at fixed $m$ in \eqref{eq:fixedmnonclt}, then
$m\to\infty$ in \eqref{eq:tailtransfer}, is the converging-together
argument.  Since $\sigma_m^2\to\|d_\theta\|_2^2=c(\theta)$, we obtain
\begin{equation}
 \Var_fM_{\theta,n}\sim c(\theta)n\log_2n,qquad
 {M_{\theta,n}-\E_fM_{\theta,n}\over\sqrt{c(\theta)n\log_2n}}
 \Longrightarrow N(0,1).                                  \label{eq:Mnonuniform}
\end{equation}
The endpoint sum has variance $O_{f,\theta}(n)$: truncate its logarithmic
singularity, use cell stabilization for the bounded part, and let the
$L^2$ tail tend to zero.  Together with \eqref{eq:VarLnon}, this proves:

\begin{theorem}[Nonuniform strict-stopped LISTER theorem]\label{thm:nonunifV}
Under \eqref{eq:fhyp}, for fixed nonresonant $\theta$,
\begin{equation}
 \Var_fV_{\theta,n}\sim c(\theta)n\log_2n,\qquad
 {V_{\theta,n}-\E_fV_{\theta,n}\over\sqrt{c(\theta)n\log_2n}}
 \Longrightarrow N(0,1),                                  \label{eq:Vnonuniform}
\end{equation}
and the same limit holds with denominator $\sqrt{\Var_fV_{\theta,n}}$.
\end{theorem}
The universality of $c(\theta)$ is now transparent: local conditional
centering changes bracket means by a summable physical-scale error, whereas
the $n\log n$ active-level mass multiplies the Lebesgue energy of the
canonical innovation.

\clearpage
\section{Nonuniform adaptive discrepancy by stabilization}\label{sec:stabilization}
The finite-tree identity \eqref{eq:finite-tree} is distribution-free.  We
write $D(Y)$ for a barred root defect and begin with a local profile $h$
whose density ratio is bounded as in \eqref{eq:flat}.  If $Y_1,\ldots,Y_n,X$
are independent $h$-points, define the root add-one increment
\begin{equation}Z_n=D(Y_1,\ldots,Y_n,X)-D(Y_1,\ldots,Y_n). \label{eq:Zn}
\end{equation}

\begin{lemma}[Root stabilization]\label{lem:rootstab}
Uniformly over physical-cell profiles and every $0<\alpha<1$,
\begin{equation}
 \Pp(Z_n\ne0)\le {C\over n+1},\qquad
 \|Z_n\|_2\le C_\alpha(n+1)^{-\alpha/2}.                 \label{eq:rootstab}
\end{equation}
\end{lemma}
\begin{proof}
Let $I^{*j}$ be the $j$-fold middle half of the normalized cell and
$S_j=\{X\in I^{*j}\}$.  Repeated application of the coefficient-one inward
copy, before making any conditioning assertion, gives for every $L$
\begin{equation}
 \{Z_n\ne0\}\subset S_L\cup
 \bigcup_{0\le j<L}(S_j\cap G_j^c),                      \label{eq:supportinc}
\end{equation}
where $G_j$ says that the fixed old-sample guard at level $j$ is filled.
Density-ratio bounds and the finite number of guard atoms give
\begin{equation}
 \Pp(S_j)\le C2^{-j},\qquad \Pp(G_j^c)\le Ce^{-cn2^{-j}}. \label{eq:twoprobs}
\end{equation}
The first event depends only on the inserted point and the second only on
the old sample.  They are therefore independent.  This is not a statement
about the law of an inward sample after conditioning on a filled guard.
Thus, on letting $L\to\infty$,
\begin{equation}
 \Pp(Z_n\ne0)\le C\sum_{j\ge0}2^{-j}e^{-cn2^{-j}}
 ={C\over n}\sum_{j\ge0}(n2^{-j})e^{-cn2^{-j}}\le {C'\over n+1}.
                                                               \label{eq:dyadicsum}
\end{equation}
For the last inequality split where $n2^{-j}\le1$; that tail is geometric,
and group the other terms in dyadic blocks, obtaining
$\sum_{k\ge0}2^ke^{-c2^k}<\infty$.

The guarded recurrence also supplies, for each finite $p$,
$\sup_{h,n}\E|D(Y_1,\ldots,Y_n)|^p<\infty$.  Consequently the same holds
for $Z_n$.  H\"older's inequality applied to
$|Z_n|^2\mathbf1_{\{Z_n\ne0\}}$, with $p>2/(1-\alpha)$, proves the second
bound in \eqref{eq:rootstab}.
\end{proof}

Let $A^+(x;Y)$ be the change in the full discrepancy after inserting $x$.
Along the physical ancestor chain $I_r(x)$, let $K_r$ be the number of old
points and let $\delta_r$ be the change of the corresponding barred root
defect.  The pathwise tree identity gives
\begin{equation}A^+(x;Y)=\sum_{r\ge0}\delta_r.             \label{eq:addonesum}
\end{equation}
Conditional on $I_r$, on $x\in I_r$, and on $K_r=k$, the old affine sample
is iid with density $f_{I_r}$ and the inserted affine point is independent
with the same density.  This legitimate cell-count conditioning permits
Lemma~\ref{lem:rootstab}.  Since $p_{I_r}\asymp_f2^{-r}$, binomial
negative-moment bounds, with $\lambda_r=m2^{-r}$, give in the dense regime
\begin{equation}\|\delta_r\|_2\le C_\alpha\lambda_r^{-\alpha/2}
 \qquad(\lambda_r\ge1).                                    \label{eq:dense}
\end{equation}
In the sparse regime, $\delta_r=0$ when $K_r=0$: both the old empty sample
and new singleton have zero barred defect.  Hence
$\Pp(K_r\ge1)\le C\lambda_r$, and the uniform high moments plus H\"older
give
\begin{equation}\|\delta_r\|_2\le C_\alpha\lambda_r^{\alpha/2}
 \qquad(\lambda_r\le1).                                    \label{eq:sparse}
\end{equation}
Both sides of the index $r\simeq\log_2m$ are geometric, so
\begin{equation}
 \sum_{r\ge0}\|\delta_r\|_2\le
 C_\alpha\sum_{r\ge0}\min\{\lambda_r,\lambda_r^{-1}\}^{\alpha/2}
 \le C_\alpha.                                             \label{eq:minksum}
\end{equation}
Apply Minkowski to finite partial sums in \eqref{eq:addonesum}; almost-sure
isolation makes the sum pathwise finite, and \eqref{eq:minksum} then gives
\begin{equation}
 \sup_{m\ge0}\E_f|A^+(X_{m+1};X_1,\ldots,X_m)|^2<\infty. \label{eq:addoneL2}
\end{equation}
No independence among depths and no covariance sign has been assumed.

Replacing one coordinate is one removal followed by one insertion.  By
exchangeability, the Efron--Stein inequality and \eqref{eq:addoneL2} yield
our distribution-free-scale comparison theorem.
\begin{theorem}[Adaptive discrepancy]\label{thm:nonunifdelta}
Under \eqref{eq:fhyp}, for every fixed $0<\theta<1$,
\begin{equation}\Var_f\Delta_{\theta,n}=O_{f,\theta}(n).   \label{eq:nonunifdelta}
\end{equation}
\end{theorem}
An alternative profile-renewal approach reduces the proof to regularity of
the difference of the two sibling mean slopes at a nonfair binomial split.
Ordinary within-profile regularity does not prove that estimate.  The
stabilization argument above makes the sibling estimate unnecessary.

\clearpage
\section{Completion of the nonuniform theorem and the mean}\label{sec:nonuniformfull}
Write $T=V+\Delta$.  Theorems~\ref{thm:nonunifV} and
\ref{thm:nonunifdelta} imply
\begin{equation}
 |\Cov_f(V,\Delta)|\le\sqrt{\Var_fV\Var_f\Delta}
 =O_{f,\theta}(n\sqrt{\log n})=o(n\log n).                \label{eq:noncov}
\end{equation}
They also imply
\begin{equation}{\Delta-\E_f\Delta\over\sqrt{n\log n}}
 \longrightarrow0\quad\hbox{in }L^2.                       \label{eq:nonDsmall}
\end{equation}
The variance identity, \eqref{eq:noncov}, and Slutsky prove
\eqref{eq:genericvar}--\eqref{eq:genericclt}; the variance asymptotic itself
converts canonical normalization to true-variance normalization.

It remains to verify the mean.  From \eqref{eq:stopped-decomp},
\eqref{eq:Fmean}, and \eqref{eq:EL},
\begin{equation}
 \E_fV_{\theta,n}={2\over\theta}\E_fL_n+E_fM_{\theta,n}
 +\E_fH_{\theta,n}={2\over\theta}n\log_2n+O_{f,\theta}(n).
                                                               \label{eq:EVmean}
\end{equation}
Here no cancellation of the last two expectations is asserted: the
conditional-bit drift and the two stopped endpoints each have absolute
first moment $O(n)$.  The finite-tree root defects have uniformly bounded
first moments and their occupied-node sum has expectation $O_{f,\theta}(n)$;
equivalently, integrate the dense/sparse add-one bounds along insertion to
obtain
\begin{equation}|\E_f\Delta_{\theta,n}|=O_{f,\theta}(n).    \label{eq:EDmean}
\end{equation}
Equations \eqref{eq:EVmean}--\eqref{eq:EDmean} prove
\eqref{eq:genericmean}.  They do not identify a second-order constant or a
log-periodic linear correction.

\clearpage
\section{Scope at resonance}\label{sec:scope}
For uniform input, Sections~\ref{sec:resstop} and \ref{sec:fullres} give the
complete theorem: at $\theta=2/Q$ the exact coboundary
\eqref{eq:cobQ} removes the repeated $n\log n$ martingale energy; external
path, endpoint, and adaptive discrepancy terms all remain on the common
order-$n$ variance scale.  Their joint signed covariance produces the
continuous one-periodic coefficient \eqref{eq:tau}; the guarded local
lower bound \eqref{eq:lower} proves strict positivity, and the joint CLT
proves the true-variance Gaussian limit.  This explains why the
normalization changes from $\sqrt{n\log n}$ to $\sqrt n$ without destroying
the CLT.

For a nonconstant density at $\theta=2/Q$, the pathwise coboundary identity
still holds and Theorem~\ref{thm:nonunifdelta} still gives
$\Var_f\Delta=O(n)$.  Nevertheless, the phase-periodic variance and CLT for
the resonant stopped LISTER, and its order-$n$ interaction with the adaptive
correction, have not been determined.  No nonuniform resonant theorem is
claimed here.  This is the principal remaining extension.

\clearpage
\section{Discussion}
The coefficient $c(\theta)$ measures the innovation repeated at active trie
levels.  Its vanishing removes an entire logarithmic factor, rather than
merely changing a constant.  The cusp shows how that factor returns near a
resonance.  The factor-three parity difference comes from an integer
endpoint profile versus cancellation at a half-integer endpoint.

The ideal and adaptive models play different roles.  Arithmetic enters the
ideal toll through endpoint orbits.  The adaptive comparison is geometric
and works for every real parameter: its guard is finite, its inward copy has
coefficient one, and its signed covariance is retained.  No result here
identifies a uniform transition window when both $n$ and $\theta$ vary.
Analogous classifications in higher dimensions would require both a
multidimensional endpoint dynamics and a replacement for the binary-trie
Walsh filtration.

\appendix
\clearpage
\section{Binary smoothing and renewal}\label{app:smoothing}
\begin{lemma}[Smoothing]\label{lem:smoothing}
Let $B_n\sim\mathrm{Bin}(n,1/2)$.  If
$u_n=2\E u_{B_n}+r_n$, $u_0=u_1=0$, and
$|r_n|\le Cn^\rho$ for some $\rho<1$, then $u_n=O(n)$.  If
$v_n=\E v_{B_n}+e_n$ with exponentially decaying $e_n$ and bounded $v_n$,
then $|v_{n+1}-v_n|=O_\beta(n^{-\beta})$ for every $\beta<1$.
\end{lemma}
\begin{proof}
Poissonize: $\widetilde u(x)=e^{-x}\sum_nu_nx^n/n!$.  Binomial splitting
becomes $\widetilde u(x)=2\widetilde u(x/2)+\widetilde r(x)$.
Iteration gives
$\widetilde u(x)=\sum_{j\ge0}2^j\widetilde r(x/2^j)$, with the finitely many
small-$x$ boundary terms included.  Since $\rho<1$, division by $x$ makes
the tail geometric, so $\widetilde u(x)=O(x)$.  Standard coefficient
comparison for nonnegative majorants, or direct binomial induction with the
self-atoms isolated, gives $u_n=O(n)$.  For the second statement couple
$B_{n+1}=B_n+I$, multiply the first-difference recursion by
$(n+1)^\beta$, and isolate $B_n=n$; the limiting contraction is
$2^{\beta-1}<1$.  Exponential forcing completes the induction.
\end{proof}

\clearpage
\section{Trie facts and the martingale CLT}\label{app:trie}
The stopped prefixes form a prefix-free set.  Conditional on their labelled
values, the unused suffixes are independent fair strings.  This follows by
checking cylinder probabilities and then applying monotone convergence.
The identities
\begin{equation}\E L_n=nS_{n-1},\qquad \Var(L_n)=O(n)           \label{eq:triefacts}
\end{equation}
follow from the binomial split recursion.  Poisson iteration gives a
continuous one-periodic $v_L$ such that $\Var(L_n)/n\to v_L(t)$ along
$\{\log_2n\}\to t$.  The same iteration applied to joint tolls gives
continuous phase covariance functions whenever the signed forcing is
$O(n^\rho)$, $\rho<1$.

We use the following standard martingale-array form: if the predictable
bracket divided by $s_n^2$ converges in probability to one and the
conditional Lindeberg sum converges in probability to zero, then the sum
divided by $s_n$ converges to $N(0,1)$.  Proposition~\ref{prop:bracket}
verifies both hypotheses rather than assuming independence of trie levels.

\clearpage
\section{Details of the resonant adaptive argument}\label{app:resonant}
We summarize the complete parity-uniform chain.  At a half-log cut
$\ell=\lfloor\tfrac12\log_2n\rfloor+O(1)$, the exact stopped residual is
\begin{equation}R_{n,\ell}=\sum_{r<\ell}\sum_{I\in\cD_r}
 D_Q(X_I)\1_{\{N_I\ge2\}}.                                  \label{eq:residual}
\end{equation}
Condition first on labelled cell allocations.  Canonical suffix blocks are
then independent.  Uniform root-defect moments control their centered
noise.  The first-difference estimate \eqref{eq:dregular}, Efron--Stein,
and multinomial coupling control the nonzero conditional mean.  At level
$r$ the variance is $O_Q(2^r)$; Minkowski across levels, without an
independence claim, yields
\begin{equation}\Var(R_{n,\ell})\le C_Q2^\ell=O_Q(\sqrt n),
 \qquad\norm{R_{n,\ell}-\E R_{n,\ell}}_2=O_Q(n^{1/4}).          \label{eq:residualbound}
\end{equation}

For even $Q$, the switching point is half-integral rather than integral.
The guard \eqref{eq:guard} prevents it from crossing an exterior active
cell; the middle-half similarity still gives exactly one inward copy and
no toll.  The endpoint orbit of $g_Q$ is finite and supplies a finite Perron
state.  The entrance/terminal covariance recursion retains its signed root
term and has forcing $O_Q(n^{1/2+\varepsilon})$.  Binary renewal therefore
produces $\kappa_Q$.  The same argument for the discrepancy produces
$\gamma_Q$.

Conditionally on labelled half-log prefixes, block pairs are independent.
Uniform fourth moments give conditional Lyapunov for every positive-variance
Cram\'er--Wold direction; conditional Chebyshev treats a degenerate
direction.  The occupancy mean is $o_{L^2}(\sqrt n)$ and
\eqref{eq:residualbound} removes the stopped residual.  Deconditioning
bounded characteristic functions gives \eqref{eq:jointres}.

For nondegeneracy, use disjoint groups of $Q$ depth-$d$ cells with three
points in the middle cell and empty guards.  There are $\Theta(n)$ candidates
when $2^d\asymp n$.  A candidate has probability bounded below.  Replacing
one label changes at most two indicators, so Efron--Stein and Chebyshev give
linearly many eligible blocks with high probability.  Inside each block,
two strict open configurations place two points in a small dyadic child and
the third respectively inside or outside its rejection distance; the local
costs differ by one.  Guard insulation makes every exterior decision
identical.  Conditional independence of the triples therefore proves
\eqref{eq:lower}.  This construction works on both integer and half-integer
switching grids.

\clearpage
\section{Fourier normalization check}\label{app:fourier}
For $m\ne0$,
\begin{equation}\widehat{R_c-2c}(m)={\sin(2\pi mc)\over\pi m}.  \label{eq:Rhat}
\end{equation}
The positive and negative odd frequencies have equal squared modulus, so
Parseval yields exactly the factor $2$ in \eqref{eq:cFourier}.  Finally
\begin{equation}
 \sum_{u\ge1,\ u\text{ odd}}{\cos(2\pi ux)\over u^2}
 =\pi^2\left({1\over8}-{1\over2}\norm{x}_{\mathbb Z}\right), \label{eq:oddcos}
\end{equation}
which inserted into the square of \eqref{eq:Psi} proves
\eqref{eq:kernel} with its sign and factor $1/2$.

\begin{thebibliography}{9}
\bibitem{BH} J.~Barnes and P.~Hut, A hierarchical $O(N\log N)$
force-calculation algorithm, \emph{Nature} \textbf{324} (1986), 446--449.
\bibitem{JS} P.~Jacquet and W.~Szpankowski, Analytic de-Poissonization and
its applications, \emph{Theoret. Comput. Sci.} \textbf{201} (1998), 1--62.
\bibitem{Mahmoud} H.~M. Mahmoud, \emph{Evolution of Random Search Trees},
Wiley, New York, 1992.
\bibitem{Billingsley} P.~Billingsley, \emph{Probability and Measure},
third ed., Wiley, New York, 1995.
\end{thebibliography}
\end{document}
